% Add this in your preamble if not already included


\section{Research Experience}
\cvSubHeadingListStart
\cvResearchHeading
  {\textbf{\makecell[l]{Generating Patient-Specific Abnormal
Gait Motions\\ for Therapists’
Observational Skill Training}}}
  {Primary Ph.D. Research}
  {Prof. Hirokazu Kato}
  {Python (PyTorch) | Deep Learning | RVQ-VAE | Unity (C\#) | Blender Python API | FastAPI  | Transformer Models | Motion Retargeting (TriLib, KeeMapRig) | HMD}

\cvItemListStart
  \cvItem{Developed an AI system that generates clinically meaningful impaired-gait animations from structured clinical text and visualizes them as 3D virtual patients for therapist training.}
  \cvItem{Retrained the MoMask text-to-motion model on a curated gait-focused dataset and therapist-supervised recordings to capture asymmetric and pathological gait.}
  \cvItem{Built a full pipeline including motion generation, Blender-based retargeting, and real-time visualization in Unity.}
\cvItemListEnd

\cvResearchHeading
  {\textbf{\makecell[l]{Augmented Reality-Based General Task Support System Using \\Large Language Models for Context-Aware and Expertise-Adaptive Guidance}}}
  {Collaborative Research}
  {Prof. Hirokazu Kato \\ \textbf{Role:} Technical collaborator and mentor for a master's research project}
  {Python | Unity (C\#) | OpenXR | FastAPI | Large Language Models (GPT-based) | Retrieval-Augmented Reasoning | Pose Tracking | AR/MR Interaction | HMD}

\cvItemListStart
  \cvItem{Contributed to the development of an AR system powered by LLMs and retrieval-augmented reasoning to deliver context-aware guidance for complex procedural tasks.}
  \cvItem{Contributed to implementing a framework that observes user actions, retrieves verified instructions from manuals, and generates expertise-adaptive AR guidance with explainability and real-time feedback.}
  \cvItem{Supported the development of applications for scientific experiments, medical device setup, and mechanical assembly workflows.}
\cvItemListEnd

\cvSubHeadingListEnd
